\chapter{Методика}
\label{ch:methods}

    Для определения теплоёмкости можно использовать формулу:
    $$Q = cm \Delta T = C \Delta T \ ,$$
    где $c$ - удельная теплоёмкость, $m$ - масса тела, $C=cm=\frac{\delta Q}{\Delta T}$ - теплоёмкость.
    
    В реальных условиях присутствует обмен с окружающей средой, поэтому формула принимает следующий вид:
    \[
    C \Delta T = P \Delta t - \lambda (T - T_k) \Delta t \ ,
    \]
    где $P \Delta t$ - количество теплоты полученное за время $\Delta t$, $\lambda$ - коэффициент теплоотдачи, $T$ - температура тела, $T_\text{к}$ - температура окружающего воздуха (комнатная). 

    Мощность нагревателя можно получить, исходя из его параметров. Коэффициент теплоотдачи можно определить по графикам зависимости температуры от времени. Температуру можно измерить термометром.